% GENERATED FILE. Edit canonical lessons, then run npm run generate:books.
% canonical-source-hash: fnv1a64:964fabcef2d7cd04
% canonical-lessons: AR-C175-sadaqa, AR-C175-mujtama, AR-C175-diyafa, AR-C175-raghma, AR-C175-aydan

\chapter{Friendship, Society, Hospitality, Despite, Also}
\label{ch:ar-friendship-society-hospitality-despite-also}
\hlchaptermodality{175}

\begin{chapteropening}
\textbf{By the end of this chapter:} \emph{I can say friendship, society, hospitality, despite and also.}

Complete the last lesson of chapter 175: I can say friendship, society, hospitality, despite and also.
\end{chapteropening}

\section[\texorpdfstring{\ar{صداقة} (ṣadāqa)}{صداقة (ṣadāqa)}]{\ar{صداقة} (ṣadāqa) — friendship}

\label{lesson:AR-C175-sadaqa}



\begin{quote}
Before the new one: say the Arabic for so, then the Arabic for so that.
\end{quote}

\subsection*{You'll want to know: \ar{صداقة}}
\textbf{\ar{صداقة}} — \emph{ṣadāqa} — \textquotedblleft{}friendship\textquotedblright{}.

\begin{cousinweb}
\textbf{\ar{صداقة}} (\emph{ṣadāqa}) is built on the native Arabic root \emph{ṣ-d-q}, \textquotedblleft{}to be truthful, sincere\textquotedblright{} (Lane). \emph{Ṣadāqa} is friendship, sincere affection.
\end{cousinweb}

\begin{grammarlens}[title={What you've built}]
One more word for this chapter.
\end{grammarlens}

\subsection*{Your turn}
\begin{itemize}
\raggedright
  \item \emph{Say it:} \emph{ṣadāqa}
  \item \emph{Say it:} \emph{ṣadāqa}, once more
  \item \emph{Say it:} say \emph{ṣadāqa}
  \item \emph{Recall:} say the Arabic for so, then the Arabic for so that, then say \emph{ṣadāqa} again
\end{itemize}

\begin{tcolorbox}[breakable,colback=teal!4,colframe=teal!35!black,title={Before you move on}]
What does \ar{صداقة} mean? (\textquotedblleft{}Friendship\textquotedblright{}.) Say it once more.
\end{tcolorbox}
\section[\texorpdfstring{\ar{مجتمع} (mujtamaʿ)}{مجتمع (mujtamaʿ)}]{\ar{مجتمع} (mujtamaʿ) — society}

\label{lesson:AR-C175-mujtama}



\begin{quote}
Before the new one: say the Arabic for so that, then the Arabic for friendship.
\end{quote}

\subsection*{You'll want to know: \ar{مجتمع}}
\textbf{\ar{مجتمع}} — \emph{mujtamaʿ} — \textquotedblleft{}society\textquotedblright{}.

\begin{cousinweb}
\textbf{\ar{مجتمع}} (\emph{mujtamaʿ}) is built on the native Arabic root \emph{j-m-ʿ}, \textquotedblleft{}to gather\textquotedblright{} (Lane). \emph{Mujtamaʿ} is the Form VIII noun of place, \textquotedblleft{}where people gather\textquotedblright{}.
\end{cousinweb}

\begin{grammarlens}[title={What you've built}]
One more word for this chapter.
\end{grammarlens}

\subsection*{Your turn}
\begin{itemize}
\raggedright
  \item \emph{Say it:} \emph{mujtamaʿ}
  \item \emph{Say it:} \emph{mujtamaʿ}, once more
  \item \emph{Say it:} say \emph{mujtamaʿ}
  \item \emph{Recall:} say the Arabic for so that, then the Arabic for friendship, then say \emph{mujtamaʿ} again
\end{itemize}

\begin{tcolorbox}[breakable,colback=teal!4,colframe=teal!35!black,title={Before you move on}]
What does \ar{مجتمع} mean? (\textquotedblleft{}Society\textquotedblright{}.) Say it once more.
\end{tcolorbox}
\section[\texorpdfstring{\ar{ضيافة} (ḍiyāfa)}{ضيافة (ḍiyāfa)}]{\ar{ضيافة} (ḍiyāfa) — hospitality}

\label{lesson:AR-C175-diyafa}



\begin{quote}
Before the new one: say the Arabic for friendship, then the Arabic for society.
\end{quote}

\subsection*{You'll want to know: \ar{ضيافة}}
\textbf{\ar{ضيافة}} — \emph{ḍiyāfa} — \textquotedblleft{}hospitality\textquotedblright{}.

\begin{cousinweb}
\textbf{\ar{ضيافة}} (\emph{ḍiyāfa}) is built on the native Arabic root \emph{ḍ-y-f}, \textquotedblleft{}to lodge as a guest\textquotedblright{} (Lane). The word \emph{ḍiyāfa} is the receiving of guests.
\end{cousinweb}

\begin{grammarlens}[title={What you've built}]
One more word for this chapter.
\end{grammarlens}

\subsection*{Your turn}
\begin{itemize}
\raggedright
  \item \emph{Say it:} \emph{ḍiyāfa}
  \item \emph{Say it:} \emph{ḍiyāfa}, once more
  \item \emph{Say it:} say \emph{ḍiyāfa}
  \item \emph{Recall:} say the Arabic for friendship, then the Arabic for society, then say \emph{ḍiyāfa} again
\end{itemize}

\begin{tcolorbox}[breakable,colback=teal!4,colframe=teal!35!black,title={Before you move on}]
What does \ar{ضيافة} mean? (\textquotedblleft{}Hospitality\textquotedblright{}.) Say it once more.
\end{tcolorbox}
\section[\texorpdfstring{\ar{رغم} (raghma)}{رغم (raghma)}]{\ar{رغم} (raghma) — despite, in spite of}

\label{lesson:AR-C175-raghma}



\begin{quote}
Before the new one: say the Arabic for society, then the Arabic for hospitality.
\end{quote}

\subsection*{You'll want to know: \ar{رغم}}
\textbf{\ar{رغم}} — \emph{raghma} — \textquotedblleft{}despite, in spite of\textquotedblright{}.

\begin{cousinweb}
\textbf{\ar{رغم}} (\emph{raghma}) is built on the native Arabic root \emph{r-gh-m}, \textquotedblleft{}dust; to be abased\textquotedblright{} (Lane). Lane's \emph{raghm} is the nose cleaving to the dust, abasement; \emph{raghma}, \textquotedblleft{}despite\textquotedblright{}, grew from it.
\end{cousinweb}

\begin{grammarlens}[title={What you've built}]
One more word for this chapter.
\end{grammarlens}

\subsection*{Your turn}
\begin{itemize}
\raggedright
  \item \emph{Say it:} \emph{raghma}
  \item \emph{Say it:} \emph{raghma}, once more
  \item \emph{Say it:} say \emph{raghma}
  \item \emph{Recall:} say the Arabic for society, then the Arabic for hospitality, then say \emph{raghma} again
\end{itemize}

\begin{tcolorbox}[breakable,colback=teal!4,colframe=teal!35!black,title={Before you move on}]
What does \ar{رغم} mean? (\textquotedblleft{}Despite, in spite of\textquotedblright{}.) Say it once more.
\end{tcolorbox}
\section[\texorpdfstring{\ar{أيضا} (ayḍan)}{أيضا (ayḍan)}]{\ar{أيضا} (ayḍan) — also, too}

\label{lesson:AR-C175-aydan}



\begin{quote}
Before the new one: say the Arabic for hospitality, then the Arabic for despite.
\end{quote}

\subsection*{You'll want to know: \ar{أيضا}}
\textbf{\ar{أيضا}} — \emph{ayḍan} — \textquotedblleft{}also, too\textquotedblright{}.

\begin{cousinweb}
\textbf{\ar{أيضا}} (\emph{ayḍan}) is built on the native Arabic root \emph{ʾ-y-ḍ}, \textquotedblleft{}to return\textquotedblright{} (Lane). \emph{Ayḍan} is \emph{ayḍ}, \textquotedblleft{}returning\textquotedblright{}, with the adverb ending \emph{-an}: \textquotedblleft{}once more\textquotedblright{}, hence \textquotedblleft{}also\textquotedblright{}.
\end{cousinweb}

\begin{grammarlens}[title={What you've built}]
One more word for this chapter.
\end{grammarlens}

\subsection*{Your turn}
\begin{itemize}
\raggedright
  \item \emph{Say it:} \emph{ayḍan}
  \item \emph{Say it:} \emph{ayḍan}, once more
  \item \emph{Say it:} say \emph{ayḍan}
  \item \emph{Recall:} say the Arabic for hospitality, then the Arabic for despite, then say \emph{ayḍan} again
\end{itemize}

\begin{tcolorbox}[breakable,colback=teal!4,colframe=teal!35!black,title={Before you move on}]
What does \ar{أيضا} mean? (\textquotedblleft{}Also, too\textquotedblright{}.) Say it once more.
\end{tcolorbox}
